\section{The reconstructed critical correspondence}
\label{sec:critical-correspondence-v143}

This section identifies the algebraic object joining the finite inverse
theorem to the real likelihood result.  The reconstruction step is a
functorial consequence of recovering the pencil.  The additional
content is a coordinate-free data construction and a relative critical
algebra whose real points split into globally labelled sections.

\subsection{The universal score scheme}
For matrix formulas we regard quadrics as symmetric bilinear forms
on a vector space; this amounts to dualizing the vector-space convention
used earlier.  Write \(\mathcal W=\Sym^2V^*\).  For a regular
pencil \(R\subset\mathcal W\), let
\(R^\times=\{K\in R:\det K\ne0\}\).
A data matrix \(S\) belongs to the same space \(\mathcal W\).
Define a section of the trivial bundle with fibre \(R^*\) on
\(R^\times\times\mathcal W\) by
\begin{equation}\label{eq:universal-score-v143}
 \eta_{K,S}(H)=-\operatorname{tr}(K^{-1}H)
       +\operatorname{tr}(S K^{-1}H K^{-1}),\qquad H\in R.
\end{equation}
Its zero scheme \(\mathfrak C_R\) is the reciprocal critical
correspondence, with its projection to the data space \(\mathcal W\).
Over the reals these are the critical equations of
\(-\log|\det K|-\operatorname{tr}(SK^{-1})\); the scheme itself
uses no choice of logarithm and is defined over the complex field.

\begin{theorem}[Reconstruction and base change of the critical correspondence]
\label{thm:critical-correspondence-v143}
Let \(n\ge3\) and \(d=n^2+2n-4\).  For every regular pencil
\(R\), the abstract scheme \(Z_R^{[d]}\) determines the
isomorphism class of the correspondence
\[
                  \mathfrak C_R\longrightarrow\mathcal W
\]
up to simultaneous congruence on the pencil and data space.
On the regular pencil parameter scheme the score sections define a
universal zero scheme.  Its pullback along any complex base change
is the corresponding relative critical scheme.  The classifying map
of a family in the framed first-relation locus therefore recovers
this relative correspondence, including over nonreduced bases.
No chosen data matrix or real form is part of the unmarked conclusion.
\end{theorem}
\begin{proof}
If \(g\) is invertible, put
\(K'=g^{\mathsf t}Kg\), \(H'=g^{\mathsf t}Hg\), and
\(S'=g^{\mathsf t}Sg\).  Cyclic invariance of trace gives
\[
                  \eta_{K',S'}(H')=\eta_{K,S}(H).
\]
Thus simultaneous congruence identifies the score sections and their
zero ideals, not just their closed points.  By
Theorem~\ref{thm:sharp-finite-pencil}, an isomorphism of the finite
neighbourhoods recovers the pencil up to projective congruence.
A linear lift of that projective transformation supplies the displayed
isomorphism of correspondences.  Different lifts differ by a common
scalar change of coordinates on the pencil and data space, so they
leave its isomorphism class unchanged.

On \(G^{\mathrm{reg}}\subset\Gr(2,\mathcal W)\), let
\(\mathcal R\) be the tautological pencil bundle.  On the determinant
open in \(\Tot\mathcal R\), inversion of the universal form is a
regular operation.  Formula~\eqref{eq:universal-score-v143} is a
section of the pullback of \(\mathcal R^*\) after taking the product
with \(\mathcal W\).  In a frame it has two entries.  Their generated
ideal is independent of the frame and its quotient commutes with
arbitrary base change, as does localization at the determinant.
This proves the relative assertion directly, without an assertion
that these score schemes are flat everywhere.  Finally the closed
immersion of Theorem~\ref{thm:finite-parameter-immersion} identifies
a framed first-relation family with its classifying morphism to
\(G\).  Pulling back the universal score construction proves the
last assertion.  It does not identify unmarked moduli stacks.
\end{proof}

\subsection{Spectral projectors and a congruence-covariant data section}
Suppose now that \(A\) is real symmetric positive definite, \(B\)
is real symmetric, and \(R=\langle A,B\rangle\) is a pencil.
Set \(C=A^{-1}B\).  Its distinct eigenvalues and multiplicities are
\[
        \alpha_1<\cdots<\alpha_r,\qquad m_1,\ldots,m_r,
        \qquad \sum_i m_i=n,
\]
where \(r\ge2\).  Its spectral idempotents are
\begin{equation}\label{eq:spectral-projectors-v143}
 P_i=\prod_{j\ne i}\frac{C-\alpha_j I}{\alpha_i-\alpha_j},
 \qquad P_iP_j=\delta_{ij}P_i,\qquad \sum_iP_i=I.
\end{equation}
They are selfadjoint for \(A\).  Define the grouped data map
\begin{equation}\label{eq:group-trace-v143}
 \tau_C:\mathcal W_{\mathbb R}\longrightarrow\mathbb R^r,
 \qquad \tau_C(S)_i=\operatorname{tr}(P_iA^{-1}S).
\end{equation}
For \(K=xB+yA\), matrix functional calculus and cyclic trace give
\begin{equation}\label{eq:matrix-grouped-likelihood-v143}
 -\log|\det K|-\operatorname{tr}(SK^{-1})
 =-\log\det A-\sum_i m_i\log|\alpha_i x+y|
                    -\sum_i\frac{\tau_C(S)_i}{\alpha_i x+y}.
\end{equation}
Thus the entire score scheme, not only its critical count, depends
on \(S\) through the surjective linear map \(\tau_C\).

Choose \(c>0\) and
\(\alpha_r<\beta_1<\cdots<\beta_{r-1}\), and set
\begin{equation}\label{eq:residue-data-section-v143}
 \begin{split}
 D(t)&=\prod_i(t-\alpha_i),\qquad Q(t)=c\prod_j(t-\beta_j),\\
 \sigma_i&=\frac{Q(\alpha_i)}{D'(\alpha_i)},\qquad
 S_0=A\sum_i\frac{\sigma_i}{m_i}P_i.
 \end{split}
\end{equation}
The form \(S_0\) is symmetric and \(\tau_C(S_0)=\sigma\), since
\(\operatorname{tr}P_i=m_i\).  More generally the assignment
\(\sigma\mapsto A\sum_i(\sigma_i/m_i)P_i\) is a right inverse
of \(\tau_C\); all data with these grouped values are
\(S_0+\ker\tau_C\).  Under simultaneous congruence,
\(C\mapsto g^{-1}Cg\), \(P_i\mapsto g^{-1}P_i g\), and
\(S_0\mapsto g^{\mathsf t}S_0g\).  No eigenbasis has been chosen.

\begin{theorem}[A totally real finite etale critical cover]
\label{thm:real-critical-cover-v143}
Fix the multiplicities \(m_1,\ldots,m_r\).  On the locus of supplied
real definite pairs \((A,B)\) with these ordered multiplicities,
there is a nonempty fibrewise open semialgebraic set of data defined
by the following condition: the polynomial
\begin{equation}\label{eq:interpolation-polynomial-v143}
 Q_S(t)=\sum_i\tau_C(S)_i\frac{D(t)}{t-\alpha_i}
\end{equation}
has positive leading coefficient and \(r-1\) distinct real roots,
all greater than \(\alpha_r\).
On this set the reciprocal critical scheme has exactly \(2r-3\)
points, all real and nondegenerate.  They form globally labelled
real analytic semialgebraic sections: \(r-1\) have their spectral
coordinate in \((\alpha_i,\alpha_{i+1})\), and \(r-2\) in the
gaps between successive roots of \(Q_S\).

In algebraic spectral coordinates the same family is the restriction
of a finite etale algebra of rank \(2r-3\) over a specified Zariski
open base.  Its construction commutes with arbitrary base change on
that base.  The formula in \eqref{eq:residue-data-section-v143}
provides a congruence-covariant section of the data map, and its
entire affine fibre has the same critical scheme.
\end{theorem}
\begin{proof}
Write \(Q=Q_S\), and put
\begin{equation}\label{eq:relative-score-polynomial-v143}
 H(t)=\sum_i(n-m_i)\frac{D(t)}{t-\alpha_i},\qquad
 F(t)=nD(t)Q'(t)-H(t)Q(t).
\end{equation}
Both apparent leading terms have degree \(2r-2\) and cancel.
Writing \(Q=c\prod_j(t-\beta_j)\), the coefficient of
\(t^{2r-3}\) is
\begin{equation}\label{eq:leading-score-coefficient-v143}
 a=c\left(n\sum_j\beta_j-\sum_i(n-m_i)\alpha_i\right)>0.
\end{equation}
The strict inequality follows from \(\beta_j>\alpha_r\),
\(n-m_i>0\), and \(\sum_i(n-m_i)=n(r-1)\).
Hence \(F\) has degree exactly \(2r-3\).

At \(\alpha_i\) and \(\beta_j\) one has
\[
 F(\alpha_i)=-(n-m_i)D'(\alpha_i)Q(\alpha_i),\qquad
 F(\beta_j)=nD(\beta_j)Q'(\beta_j).
\]
The signs alternate along each string.  The first string gives a
root in every \((\alpha_i,\alpha_{i+1})\), and the second gives
one in every \((\beta_j,\beta_{j+1})\).  These \(2r-3\)
disjoint intervals exhaust the degree.  All roots are simple and
\(F\) is coprime to \(DQ\).  This also treats \(r=2\), when
the second string of intervals is empty.

We next prove the scheme statement and exclude the missing coordinate
chart uniformly.  Take the polynomial ring over \(\mathbb Q\) in
\(\alpha_1,\ldots,\alpha_r,\beta_1,\ldots,\beta_{r-1},c\),
and localize at \(c\), the products of distinct spectral and
\(\beta\)-differences, \(a\), \(\disc_tF\), and
\(\operatorname{Res}_t(F,DQ)\).  Call this ring \(\mathscr B\).
Every point of the separated real chamber lies in
\(\Spec\mathscr B(\mathbb R)\), by the preceding sign argument.
For a linear \(F\), its discriminant is understood as \(1\).

The partial fraction identity \(Q/D=\sum_i\sigma_i/(t-\alpha_i)\)
shows that
\[
 \sum_i\sigma_i=c,\qquad
 \sum_i\alpha_i\sigma_i
       =c\left(\sum_i\alpha_i-\sum_j\beta_j\right).
\]
On the chart \(x=0\) of the invertible-pencil locus the score
in \(y\) forces \(y=c/n\).  Substitution in the score in
\(x\) then forces \(a=0\).  Since \(a\) is inverted, the
quotient of the critical coordinate ring by \((x)\) is zero.
Thus \(x\) is a unit on the whole critical scheme, including
nilpotent test rings; excluding this chart is not merely a pointwise
argument.

Set \(t=-y/x\) and
\(U(t)=\sum_i\sigma_i/(\alpha_i-t)=-Q(t)/D(t)\).
In these coordinates the grouped function is
\(-n\log x-\sum_i m_i\log(\alpha_i-t)-U(t)/x\),
with the logarithms used only to write its rational differential.
The score in \(x\) gives \(x=U/n\); the remaining score becomes
\(-F/(DQ)\).  Consequently the critical coordinate algebra is
\begin{equation}\label{eq:finite-etale-score-algebra-v143}
 \mathscr B[t]/(F/a),\qquad
 x=-\frac{Q(t)}{nD(t)},\qquad
 y=\frac{tQ(t)}{nD(t)}.
\end{equation}
The resultant makes \(D\) and \(Q\) units in this algebra, so
these formulas and the coordinate change are genuine inverse
homomorphisms.  The polynomial \(F/a\) is monic of degree
\(2r-3\); its discriminant is a unit.  The algebra is therefore
finite free of that rank, and its derivative is a unit modulo the
polynomial, which proves that it is etale.  These algebraic
identities persist under every base change of \(\mathscr B\).
The statement is local in spectral coordinates; no global flatness
of the unrestricted critical correspondence is inferred.

At a critical point the second derivative in the eliminated \(x\)
direction is \(-n/x^2\).  The complementary Schur factor is
\(-F'(t)/(D(t)Q(t))\).  Both are nonzero, so the Hessian is
nondegenerate.  The change from \((x,t)\) to \((x,y)\) has
Jacobian determinant \(-x\ne0\), preserving this conclusion.
Each specified real interval contains a unique root.  The implicit
function theorem therefore supplies a real analytic root function
locally, and uniqueness glues these functions globally with the
stated labels.  Their graphs are semialgebraic, since they are
specified by a polynomial equation and the interval inequalities.
The ordered roots \(\beta_j\) of \(Q_S\) have the same local
analytic and global semialgebraic property.

For fixed \(\alpha\), the map from the \(r\) coefficients of
\(Q\) to \(\sigma\) is a linear isomorphism, with inverse
\eqref{eq:interpolation-polynomial-v143}.  Polynomials with the
specified simple separated roots and positive leading coefficient
form a nonempty open semialgebraic coefficient set.  Surjectivity
and the explicit right inverse of \(\tau_C\) make its preimage
in the full data space nonempty and open.  Spectral projectors give
the same formulas along any real analytic family in the fixed
multiplicity locus.  Their congruence covariance was verified above;
the kernel directions of \(\tau_C\) do not occur in the score
by \eqref{eq:matrix-grouped-likelihood-v143}.  This completes all
the assertions.
\end{proof}

\begin{corollary}[The finite-failure-to-critical construction]
\label{cor:failure-to-critical-v143}
An order-\(d\) finite failure invariant of a regular complex pencil,
together with a supplied real definite realization and the spectral
choices in \eqref{eq:residue-data-section-v143}, determines the
above finite etale critical family up to simultaneous congruence.
The all-real critical count and labelled real sections are independent
of an eigenbasis.  The real choices are additional inputs, not
outputs of complex unmarked reconstruction.
\end{corollary}
\begin{proof}
Apply Theorem~\ref{thm:critical-correspondence-v143}, use the
spectral sheaf to read the eigenvalues and their multiplicities in
the supplied definite realization, and apply
Theorem~\ref{thm:real-critical-cover-v143}.  Equations
\eqref{eq:spectral-projectors-v143}--\eqref{eq:residue-data-section-v143}
verify invariance under all changes of eigenbasis and simultaneous
congruence.  The proof of total reality is an additional spectral
argument, not a premise of the inverse theorem.
\end{proof}
